\section{Related work}

\subsection{Scaling Up Representation Learning for Natural Language}
Learning representations of natural language has been shown to be useful for a wide range of NLP tasks and has been widely adopted~\citep{mikolov2013distributed,le-mikolov:2014,dai2015semi, peters2018deep,devlin2018bert,radford2018improving,radford2019language}.
One of the most significant changes in the last two years is the shift from pre-training word embeddings, whether standard~\citep{mikolov2013distributed,pennington-etal-2014-glove} or contextualized~\citep{NIPS2017_7209,peters2018deep}, to full-network pre-training followed by task-specific fine-tuning~\citep{dai2015semi, radford2018improving,devlin2018bert}.
In this line of work, it is often shown that larger model size improves performance.
For example, \cite{devlin2018bert} show that across three selected natural language understanding tasks, using larger hidden size, more hidden layers, and more attention heads always leads to better performance.
However, they stop at a hidden size of 1024, presumably because of the model size and computation cost problems. 
% We show that, under the same setting, increasing the hidden size to 2048 leads to model degradation and hence worse performance.
% Therefore, scaling up representation learning for natural language is not as easy as simply increasing model size.

It is difficult to experiment with large models due to computational constraints, especially in terms of 
GPU/TPU memory limitations.
Given that current state-of-the-art models often have hundreds of millions or even billions of parameters, we can easily hit  memory limits.
To address this issue, \cite{chen2016training} propose a method called gradient checkpointing to reduce the memory requirement to be sublinear at the cost of an extra forward pass. 
\cite{gomez2017reversible} propose a way to reconstruct each layer's activations from the next layer so that they do not need to store the intermediate activations.
Both methods reduce the memory consumption at the cost of speed. \cite{raffel2019exploring} proposed to use model parallelization to train a giant model.
In contrast, our parameter-reduction techniques reduce memory consumption and increase training speed. 


\subsection{Cross-layer parameter sharing}
The idea of sharing parameters across layers 
%cross-layer parameter sharing 
has been previously explored with the Transformer architecture~\citep{vaswani2017attention}, but this prior work has focused on training for standard encoder-decoder tasks rather than the pretraining/finetuning setting. Different from our observations, \cite{dehghani2018universal} show that  networks with cross-layer parameter sharing (Universal Transformer, UT) get better performance on language modeling and subject-verb agreement than the standard transformer. Very recently, \cite{bai2019deep} propose a Deep Equilibrium Model (DQE) for transformer networks and show that DQE can reach an equilibrium point for which the input embedding and the output embedding of a certain layer stay the same. Our observations show that our embeddings are oscillating rather than converging. \cite{Hao_2019} combine a parameter-sharing transformer with the standard one, which further increases the number of parameters of the standard transformer.  

\vspace{-5pt}
\subsection{Sentence Ordering Objectives}
ALBERT uses a pretraining loss based on predicting the ordering of two consecutive segments of text. Several researchers have experimented with pretraining objectives that similarly relate to discourse coherence. 
Coherence and cohesion in discourse have been widely studied and many phenomena have been identified that connect neighboring text segments~\citep{DBLP:journals/cogsci/Hobbs79,halliday-76,grosz-etal-1995-centering}. 
Most objectives found effective in practice are quite simple. 
Skip-thought~\citep{Kiros2015skipthought} and FastSent~\citep{hill2016fastsent} sentence embeddings are learned by using an encoding of a sentence to predict words in neighboring sentences. 
Other objectives for sentence embedding learning include predicting future sentences rather than only neighbors~\citep{gan-etal-2017-learning} and predicting explicit discourse markers~\citep{jernite2017discourse,nie-etal-2019-dissent}. 
Our loss is most similar to the sentence ordering objective of \citet{jernite2017discourse}, where sentence embeddings are learned in order to determine the ordering of two consecutive sentences. 
Unlike most of the above work, however, our loss is defined on textual segments rather than sentences. BERT~\citep{devlin2018bert} uses a loss based on predicting whether the second segment in a pair has been swapped with a segment from another document.
We compare to this loss in our experiments and find that sentence ordering is a more challenging pretraining task and more useful for certain downstream tasks.
Concurrently to our work, \cite{wang2019structbert} also try to predict the order of two consecutive segments of text, but they combine it with the original next sentence prediction in a three-way classification task rather than empirically comparing the two. 
